% Introduction: the objects of the paper. This chapter is unnumbered (\chapter*), so its items
% are numbered 0.1, 0.2, ... and the paper's own numbering starts in the next chapter.
\chapter*{Introduction}
\addcontentsline{toc}{chapter}{Introduction}

Throughout this paper, $a$ denotes an integer such that $-a$ is not a square in $\QQ$,
$f := X^2 + a \in \ZZ[X]$, and $f_0 := X$, $f_{n+1} := f(f_n) = f_n^2 + a$ for all $n \ge 0$, are
the iterates of $f$. Let $c_1 := -a$ and $c_{n+1} := f(c_n) = c_n^2 + a$
$(= f_n(-a) = f_n(a) = f_{n+1}(0))$ for $n \ge 1$. There exists an integer sequence $(b_n)_{n
\ge 1}$ such that for all $n \ge 1$ we have $c_n = \prod_{d \dvd n} b_d$
(Lemma~\ref{lem:1.1}). Let $K_n$ be the splitting field of $f_n$ over $\QQ$ and denote by
$\Om_n := \Gal(K_n/\QQ) = \Gal(f_n/\QQ)$ its Galois group over the rational numbers.

This paper deals with the problem of determining $\Om_n$. Following R.~W.~K.~Odoni's paper
\cite{odoni}, we generalize his result concerning the case $a = 1$ to get for all $n \ge 1$:
\[ \Om_n \cong \Ct{n}, \text{ if and only if } b_1, b_2, \dots, b_n
   \text{ are $2$-independent in } \QQ \tag{$*$} \]
(Theorem~\ref{thm:section1}). $\Om_n$ always injects into $\Ct{n}$, the $n$-fold wreath product
of $C_2$ with itself. In order to prove the condition on the right hand side of $(*)$ (at least
for some $a \in \ZZ$), we consider more generally iteration sequences associated to certain even
polynomials $g \in \ZZ[X^2]$ (Chapter~2). Specializing these results, we get with $(*)$: if
($a > 0$ and $a \equiv 1$ or $2 \bmod 4$) or ($a < 0$ and $a \equiv 0 \bmod 4$), then
$\Om_n \cong \Ct{n}$ for all $n \ge 1$ (Theorem~\ref{thm:section3}). Moreover, this gives a proof
of the conjecture of John~E.~Cremona \cite{cremona}, who verified numerically that, for $a = 1$,
$|b_n|$ is not a square when $2 \le n \le 5 \cdot 10^7$.

\begin{definition}[The iterates $f_n$]
  \label{def:iterates}
  \lean{QuadraticIterates.iteratedPoly, QuadraticIterates.iteratedPoly_zero,
    QuadraticIterates.iteratedPoly_succ, QuadraticIterates.iteratedPoly_succ_comp,
    QuadraticIterates.iteratedPoly_add, QuadraticIterates.map_iteratedPoly,
    QuadraticIterates.monic_iteratedPoly, QuadraticIterates.natDegree_iteratedPoly,
    QuadraticIterates.ncard_rootSet_iteratedPoly_le}
  \leanok
  $f_0 := X$ and $f_{n+1} := f_n^2 + a$. These are monic of degree $2^n$ --- so $f_n$ has at
  most $2^n$ roots --- and they satisfy
  $f_{m+n} = f_m \circ f_n$; in particular $f_{n+1} = f_n \circ (X^2 + a)$, which is the form in
  which the iterate is taken apart in Chapter~1.
  \formalnote{\texttt{iteratedPoly} is defined over an arbitrary commutative semiring; the
  sequence of the paper is the case $R = \ZZ$. Since forming iterates commutes with ring
  homomorphisms (\lname{map\_iteratedPoly}), the image of $f_n$ in $\QQ[X]$ --- written
  \texttt{f}$\QQ$\texttt{[a, n]} in Lean --- is again an iterate, and all statements over $\QQ$
  are about that polynomial.}
\end{definition}

\begin{definition}[The fields $K_n$ and the groups $\Om_n$]
  \label{def:galois}
  \uses{def:iterates}
  \lean{QuadraticIterates.splittingField, QuadraticIterates.GaloisGroup,
    QuadraticIterates.splittingField.instIsSplittingField,
    QuadraticIterates.splittingField.instFiniteDimensional,
    QuadraticIterates.splittingField.instIsGalois,
    QuadraticIterates.splittingField_zero_eq_bot,
    QuadraticIterates.card_galoisGroup_eq_finrank}
  \leanok
  $K_n \subseteq \Qbar$ is the subfield generated by the roots of $f_n$, and
  $\Om_n = \Gal(f_n/\QQ)$. One has $K_0 = \QQ$ and $\#\Om_n = [K_n : \QQ]$.
  \formalnote{$K_n$ is realized concretely as an \texttt{IntermediateField} of
  \texttt{AlgebraicClosure} $\QQ$ --- namely \texttt{IntermediateField.adjoin} of the root set ---
  rather than as an abstract splitting field, so that the tower $K_n \subseteq K_{n+1}$ is an
  inclusion of subfields of one fixed field and relative degrees are
  \texttt{IntermediateField.relfinrank}. That this field \emph{is} a splitting field, is finite
  over $\QQ$ and is Galois are instances. $\Om_n$ is Mathlib's \texttt{Polynomial.Gal}.}
\end{definition}

\begin{definition}[The iterated wreath power $\Ct{n}$]
  \label{def:wreath}
  \lean{QuadraticIterates.WreathPower, QuadraticIterates.card_wreathPower}
  \leanok
  $\Ct{n}$ is the $n$-fold iterated regular wreath product of $C_2$ with itself; it has order
  $2^{2^n - 1}$.
  \formalnote{This is Mathlib's \texttt{IteratedWreathProduct}, applied to the group
  \texttt{Multiplicative (ZMod 2)}. Statements of the form ``$\Om_n \cong \Ct{n}$'' are
  formalized as \texttt{Nonempty} of a \texttt{MulEquiv} between \texttt{GaloisGroup a n} and
  \texttt{WreathPower n}: the \emph{existence} of an isomorphism, which is what the paper
  asserts, rather than a chosen one.}
\end{definition}

\begin{definition}[The sequence $c_n$]
  \label{def:cseq}
  \uses{def:iterates}
  \lean{QuadraticIterates.cSeq, QuadraticIterates.cSeq_zero, QuadraticIterates.cSeq_one,
    QuadraticIterates.cSeq_succ, QuadraticIterates.cSeq_two,
    QuadraticIterates.cSeq_succ_eq_neg_one_pow_mul_eval,
    QuadraticIterates.intCast_cSeq_eq_ite_mul_eval_zero}
  \leanok
  $c_1 = -a$ and $c_{n+1} = c_n^2 + a$ for $n \ge 1$. One has $c_{n+1} = (-1)^{2^n} f_n(a)$, so
  $c_{n+1} = f_n(a) = f_{n+1}(0)$ for $n \ge 1$, while $c_1 = -f_1(0)$.
  \formalnote{\texttt{cSeq} is total, with the junk value $c_0 = 0$; the recursion
  \lname{cSeq\_succ} therefore carries the hypothesis $n \ge 1$. The junk value is not
  arbitrary: it is what makes $c$ a \emph{strong divisibility sequence}
  ($\gcd(c_m, c_n) = |c_{\gcd(m,n)}|$, see Lemma~\ref{lem:1.1}), which is how integrality of the
  $b_n$ is obtained. \texttt{cSeq} is the specialization to $g = X^2 + a$, $\varepsilon = -1$ of
  the general sequence of Definition~\ref{def:gammaseq}.}
\end{definition}

\begin{definition}[The Möbius factors $b_n$]
  \label{def:bseq}
  \uses{def:cseq}
  \lean{QuadraticIterates.bSeq, QuadraticIterates.bSeq_one,
    QuadraticIterates.bSeq_eq_moebiusFactorR}
  \leanok
  $b_n := \prod_{d \dvd n} c_d^{\mu(n/d)}$, where $\mu$ is the Möbius function. In particular
  $b_1 = c_1 = -a$.
  \formalnote{The paper \emph{constructs} the $b_n$ inside Lemma~1.1~b) and proves integrality
  afterwards. The formalization instead \emph{defines} $b_n$ at once, as an element of $\ZZ$: it
  is \texttt{moebiusFactorR} of the sequence $c$, the preimage in $\ZZ$ of the rational Möbius
  product (junk when that product is not integral). Lemma~\ref{lem:1.1}~b) is then the statement
  that this preimage really maps to the rational product $\prod_{ed = n} c_d^{\mu(e)}$
  (\lname{intCast\_bSeq}), i.e. that the product is an integer. The rational product itself is
  \texttt{moebiusFactorK} with $K = \QQ$, written over the divisor antidiagonal of $n$
  (Lemma~\ref{lem:moebiusfactor}).}
\end{definition}

\begin{definition}[$2$-independence]
  \label{def:twoindep}
  \lean{QuadraticIterates.TwoIndependent, QuadraticIterates.TwoIndependent.ne_zero,
    QuadraticIterates.twoIndependent_iff_linearIndependent,
    QuadraticIterates.TwoIndependent.of_snoc, sqClass}
  \leanok
  Nonzero rational numbers $a_1, \dots, a_n$ are called \emph{$2$-independent} if their residue
  classes in the $\FF_2$-vector space $\QQ^*/(\QQ^*)^2$ are linearly independent.
  \formalnote{\texttt{TwoIndependent} is stated in the equivalent elementary form ``no nonempty
  subfamily has a square product'', which is what the proofs use, for a family of elements of any
  commutative monoid; that all $a_i$ are nonzero is then a consequence, $0$ being a square
  (\lname{TwoIndependent.ne\_zero}). The translation into linear independence of the classes
  $\texttt{sqClass}(a_i)$ in $L^*/(L^*)^2$, for any field $L$, is
  \lname{twoIndependent\_iff\_linearIndependent}. The families of the paper are indexed by
  \texttt{Fin n}, and $c_1, \dots, c_n$ appears as the family $i \mapsto c_{i+1}$.}
\end{definition}
